\begin{scanpage}{0200}{182}
\[
\psi_\pm(z)=\frac{\displaystyle\ln\frac{-1}{\gamma^2q_\pm}}{\displaystyle\ln\frac{2ix}{\gamma}-E(2q_\pm x)e^{-2iq_\pm x}},\qquad\Biggr\}\tag{28.21}
\]
\noindent where
\[
\ln(-1)=i\pi,\quad\ln i=i\frac\pi2.\tag{28.22}
\]
\[
x=\frac{kz}{q},\quad q=(ka)^2,\quad q_\pm=q\frac{1\mp\cos\vartheta}{2}\tag{28.23}
\]
\noindent and
\[
E(y)=\operatorname{Ci}y+i\operatorname{si}y=-\int_y^\infty\frac{e^{it}}t\,dt.\tag{28.24}
\]
The integral cosine $\operatorname{Ci}y$ and the integral sine $\operatorname{si}y$ are determined by the relationships
\[
\operatorname{Ci}y=-\int_y^\infty\frac{\cos t}t\,dt,\quad
\operatorname{si}y=-\int_y^\infty\frac{\sin t}t\,dt\tag{28.25}
\]
and are thoroughly tabulated functions.

The equations written above for the current in a finite conductor are distinguishable by their visualizability, and they enable one to liken the conductor to a section of a transmission line in which, however, the attenuation of the current waves takes place, not according to an exponential law, but according to a more complicated law which is determined by Equations (28.21). In addition, the diffraction character of the problem is reflected in the equations. The conductor's specific features as a diffraction object are included in the very slow attenuation of the current waves. As a consequence of this, it is impossible to limit oneself to considering only secondary and tertiary diffractions, and it is necessary to sum all the reflected waves. As a result of such a summation, a ``resonance denominator'' $D$ appears which takes into account the resonance properties of a thin, finite length conductor.
\end{scanpage}
